\section{SHAP is not feature selection --- attribution vs value of information}

A common and expensive confusion. Three different questions hide under
``feature importance'':

\begin{enumerate}
  \item \textbf{What does THIS fitted model lean on?} $\to$ SHAP
        (attribution).
  \item \textbf{What is a feature or group worth --- what does accuracy
        lose if I remove it and RETRAIN?} $\to$ ablation
        (value of information).
  \item \textbf{What is the smallest feature set that keeps accuracy?}
        $\to$ forward/backward selection, or L1.
\end{enumerate}

The answers can disagree wildly, and correlated features are the reason.

\subsection{Why they disagree: credit splitting}

SHAP divides a prediction's credit fairly across features. If ten
features carry copies of the same signal, each gets $\sim$1/10 of the
credit and each ranks low --- though the signal itself is the most
valuable thing in the data. Remove one of the ten and retrain: the model
reorganizes around the other nine and loses almost nothing. So:

\begin{itemize}
  \item high SHAP $\neq$ irreplaceable,
  \item low SHAP $\neq$ useless,
  \item and no SHAP rank tells you what a RETRAINED model would do ---
        SHAP is computed inside one frozen model.
\end{itemize}

A field example (from a practitioner): 200 candidate features,
AUROC 0.845. Top-10 by SHAP shared \emph{zero} features with the 8
found by forward selection (AUROC 0.836). One of the 8 scored 0.74
alone, yet SHAP ranked it 6th of 8. Nothing malfunctioned: under heavy
redundancy many disjoint subsets are near-optimal, the lone-strong
feature's information also arrived via correlated colleagues (so its
marginal credit was small), and greedy forward selection is itself
unstable --- rerun on a bootstrap and the chosen 8 change.

\subsection{Our own case: ``solar is price driver \#1'' --- tested, survived}

SHAP on the price LightGBM ranks the solar forecast first
(18.7 EUR mean $|$SHAP$|$) above yesterday's price (14.1). Suspicion:
our $\sim$29 correlated price columns split credit, deflating each and
inflating solar's rank. So we ran the retrain test
(\texttt{reports/backtests/2026-07-17\_price\_group\_ablation.md},
walk-forward, Jan--Jul 2026): drop the RES group, MAE +3.54 EUR/MWh ---
the LARGEST group cost, above even the whole price-lag block (+2.80).
The suspicion was reasonable and WRONG; the claim now stands on two
independent methods. Honest caveats: the test window is the
solar-heavy season (winter would shift value toward scarcity signals),
and load lags turned out to be dead weight for the price model
($-0.04$; pruning candidate).

Do not skip the general lesson because this case agreed: attribution
magnitude and retrain delta are DIFFERENT UNITS answering different
questions --- solar's 18.7 EUR mean $|$SHAP$|$ is not ``worth 18.7 EUR
of MAE'' (the retrain says 3.54 for the whole group). Only the ranking
was confirmed, and only at group level.

The load model shows the same double-check: SHAP said the TSO forecast
dominates, and the retrain ablation confirmed it (its removal costs
+1.97pp MAPE, 96\% of total skill). When both methods agree, the
conclusion is solid. When they disagree, say WHY before quoting either.

\subsection{What the original paper actually promises}

Lundberg \& Lee (2017), read closely, confirm all of this:

\begin{itemize}
  \item \textbf{Local accuracy}: the attributions are DEFINED to sum to
        one prediction minus the model's average output ---
        $f(x) = \phi_0 + \sum_i \phi_i$. A decomposition of a fitted
        model's output. The paper's stated goal is user trust and model
        understanding; feature selection is never claimed.
  \item \textbf{The estimators assume feature independence} (their
        Eq.~11). Our 29 correlated price columns violate this exactly
        where it matters --- credit assignment among correlated features
        is where the practical approximations are weakest.
  \item \textbf{The retraining variant exists}: their ``Shapley
        regression values'' retrain on every feature subset ($2^M$
        models) and average performance deltas --- value of information,
        done exhaustively. SHAP's innovation is AVOIDING retraining for
        speed; that choice is precisely what decouples it from
        predictive value. Our group ablation is the affordable cousin:
        one coalition per group, retrain, measure.
  \item \textbf{Why SHAP anyway}: the consistency theorem. Gain-based
        importances (LightGBM's built-in) can DECREASE when a feature's
        true contribution increases; SHAP provably cannot. Among
        attribution tools it is the right one --- the mistake is only
        asking attribution to do selection.
\end{itemize}

\subsection{Practice rules for this repo}

\begin{itemize}
  \item SHAP is for explaining the shipped model's daily forecast
        (hard rule 3). Never for feature selection.
  \item Any ``driver'' claim in a card or README must cite the retrain
        ablation, not the SHAP rank, or state both.
  \item Feature pruning decisions: ablation on groups, LASSO path for
        linear models, never SHAP rank alone.
\end{itemize}

Interview line: ``SHAP explains the model you have; ablation prices the
data you feed it. Correlated features make the two disagree, and the
size of the disagreement is itself a useful measurement.''
